\section{Artin rings}
\label{am-ch8-artin}
\begin{definition}[Artin ring]
\label{am-def-artin-ring}
\source{atiyah-macdonald-commutative-algebra:page-0098}
An Artin ring satisfies the descending chain (minimal) condition on ideals.
\end{definition}
\begin{proposition}[Primes in an Artin ring]
\label{am-prop-8-1}
\dcref{8.1}
\source{atiyah-macdonald-commutative-algebra:page-0098}
Every prime ideal of an Artin ring is maximal.
\end{proposition}
\begin{proof}
If $\p$ is prime, $A/\p$ is an Artinian domain.  The descending chain of
principal ideals $(x)\supseteq(x^2)\supseteq\cdots$ stabilizes for $x\ne0$;
$x^n=x^{n+1}y$ and cancellation gives $xy=1$.  Thus the quotient is a field.
\end{proof}
\begin{corollary}[Radicals in Artin rings]
\label{am-cor-8-2}
\dcref{8.2}
\source{atiyah-macdonald-commutative-algebra:page-0098}
In an Artin ring the nilradical equals the Jacobson radical.
\uses{am-prop-8-1,am-prop-1-8,am-def-jacobson}
\end{corollary}
\begin{proof}
The nilradical is the intersection of all primes, and Proposition~\ref{am-prop-8-1}
identifies these with maximal ideals.
\end{proof}
\begin{proposition}[Finitely many maximal ideals]
\label{am-prop-8-3}
\dcref{8.3}
\source{atiyah-macdonald-commutative-algebra:page-0098}
An Artin ring has only finitely many maximal ideals.
\end{proposition}
\begin{proof}
The set of finite intersections of maximal ideals has a minimal member
$\m_1\cap\cdots\cap\m_n$.  Intersecting it with any maximal $\m$ leaves it
unchanged; prime avoidance (Proposition~\ref{am-prop-1-11}) forces
$\m=\m_i$ for some $i$.
\uses{am-prop-1-11}
\end{proof}
\begin{proposition}[Nilradical is nilpotent]
\label{am-prop-8-4}
\dcref{8.4}
\source{atiyah-macdonald-commutative-algebra:page-0098}
The nilradical of an Artin ring is nilpotent.
\end{proposition}
\begin{proof}
The powers of the nilradical stabilize, say $\Nil^k=\Nil^{k+1}=I$.  If
$I\ne0$, choose a minimal ideal $\cideal$ with $I\cideal\ne0$ and $x\in\cideal$
with $xI\ne0$.  Minimality gives $\cideal=(x)=xI$; writing $x=xy$ with
$y\in I$ and iterating yields $x=xy^n=0$ for large $n$, contradiction.
\end{proof}
\begin{definition}[Dimension of a ring]
\label{am-def-ring-dimension}
\source{atiyah-macdonald-commutative-algebra:page-0099}
The dimension is the supremum of lengths of finite strict chains of prime ideals.
\end{definition}
\begin{theorem}[Artin characterization]
\label{am-thm-8-5}
\dcref{8.5}
\source{atiyah-macdonald-commutative-algebra:page-0099}
A ring is Artinian iff it is Noetherian and has dimension $0$.
\end{theorem}
\begin{proof}
If Artinian, all primes are maximal.  The finite maximal-ideal intersection
from Proposition~\ref{am-prop-8-3} has a nilpotent power zero by
Proposition~\ref{am-prop-8-4}; Corollary~\ref{am-cor-6-11} gives Noetherianity.
Conversely, a Noetherian ring has a finite primary decomposition of zero.  Its
minimal primes are maximal when dimension is zero, and the nilradical is
nilpotent by Corollary~\ref{am-cor-7-15}; the product of sufficiently high powers
of these maximal ideals is zero, so Corollary~\ref{am-cor-6-11} gives the d.c.c.
\uses{am-prop-8-3,am-prop-8-4,am-cor-6-11,am-thm-7-13,am-cor-7-15}
\end{proof}
\begin{proposition}[Local Artin rings]
\label{am-prop-8-6}
\dcref{8.6}
\source{atiyah-macdonald-commutative-algebra:page-0099}
For a Noetherian local ring $(A,\m)$, either $\m^n\ne\m^{n+1}$ for every $n$,
or $\m^n=0$ for some $n$, in which case $A$ is Artinian.
\end{proposition}
\begin{proof}
If $\m^n=\m^{n+1}$, Nakayama's lemma gives $\m^n=0$.  Every prime contains
$\m^n$, hence has radical $\m$ and equals $\m$; dimension is zero, so use
Theorem~\ref{am-thm-8-5}.
\uses{am-prop-2-6,am-thm-8-5}
\end{proof}
\begin{theorem}[Structure theorem for Artin rings]
\label{am-thm-8-7}
\dcref{8.7}
\source{atiyah-macdonald-commutative-algebra:page-0099}
Every Artin ring is, uniquely up to isomorphism, a finite product of Artin local
rings.
\end{theorem}
\begin{proof}
Let $\m_i$ be its finitely many maximal ideals.  A common power of their
product is zero; they are pairwise coprime, so the Chinese remainder theorem
gives $A\simeq\prod_iA/\m_i^k$, each factor local Artin.  Conversely, in a
finite product of Artin local rings the kernels of the projections are pairwise
coprime and their intersection is zero; the uniqueness of isolated primary
components (Corollary~\ref{am-cor-4-11}) makes the factors unique.
\uses{am-prop-8-3,am-prop-8-4,am-prop-1-10,am-cor-4-11}
\end{proof}
\begin{proposition}[Principal Artin local rings]
\label{am-prop-8-8}
\dcref{8.8}
\source{atiyah-macdonald-commutative-algebra:page-0100}
For an Artin local ring $(A,\m)$ with residue field $k$, the following are
equivalent: every ideal is principal; $\m$ is principal; and
$\dim_k(\m/\m^2)\le1$.
\end{proposition}
\begin{proof}
The first two implications are immediate.  If the dimension is zero,
Nakayama gives $\m=0$.  If it is one, Proposition~\ref{am-prop-2-8} gives
$\m=(x)$.  Since $\m$ is nilpotent by Proposition~\ref{am-prop-8-4}, any
nonzero proper ideal lies between two consecutive powers $(x^r)$ and
$(x^{r+1})$; an element in the difference is a unit times $x^r$, so the ideal
is $(x^r)$.
\uses{am-prop-2-8,am-prop-8-4}
\end{proof}
